% Includable supplement. The division-algebra theorem is supplied separately
% in phase3/number_fields/division_norm_compiler_20261106_v1.tex.
\section{Several equations versus one equation}
\label{p3sys:combination}

Multiple equations can make the degree smaller and the construction
easier to inspect.  They do not lower the minimum possible number of
variables when degree is unrestricted: over the rationals, and over
every number field, a finite conjunction can be combined into one
equation without adding a variable.  The degree and field of the
combining polynomial must nevertheless be recorded.

\begin{proposition}\label{p3sys:sos-combination}
Let $K$ be a number field with a real embedding, and let
$F_1,\ldots,F_s\in K[T,Y_1,\ldots,Y_r]$ have degree at most $D$.
The system $F_1=\cdots=F_s=0$ has exactly the same $K$-solutions
as the single equation
\[
                         P=\sum_{i=1}^sF_i^2=0.
\]
No new variable is introduced, and $\deg P\le2D$.
For a single residual one may simply take $P=F_1$.
If $K=\mathbb Q$, $P\ne0$, all residuals have integer coefficients of
absolute value at most $H$, and there are $n=r+1$ total
coordinates, then
\[
               H(P)\le s\binom{n+D}{D}H^2.
\]
Dividing the coefficients by their common integer divisor can
only improve this bound.
\end{proposition}
\begin{proof}
Apply a real embedding of $K$ to a zero of the displayed sum.
Every summand is a nonnegative real number, so all of them vanish;
injectivity of the embedding gives $F_i=0$ in $K$.  The converse
and the degree bound are immediate.  There are at most
$M=\binom{n+D}{D}$ monomials in each residual.  For a fixed
monomial in $F_i^2$, choosing one exponent determines the other,
so its coefficient is a sum of at most $M$ products, each of
absolute value at most $H^2$.  Summing over $i$ proves the bound.
\end{proof}

The real-embedding condition is weaker than total reality.
It is essential to this particular combining formula: over
$\mathbb Q(i)$, the two nonzero values $1,i$ have squares summing
to zero.  A replacement over other number fields is supplied by
division algebras.

\begin{corollary}\label{p3sys:division-combination}
For any number field $K$, a system of $s\ge1$ polynomial equations
of degrees at most $D$ can be combined into one equation in the
same coordinates of degree at most
\[
                         D\lceil\sqrt{s}\rceil.
\]
For a specified $K$, a combining polynomial with rational
coefficients suffices.  This is a statement about equivalence
of solution sets over that field; it does not identify a changed
field's solution set with its ring of integers.
\end{corollary}
\begin{proof}
Put $d=\lceil\sqrt{s}\rceil$.  The division-algebra construction
of Theorem~\ref{thm:division-norm-conjunction-phase3} supplies a
homogeneous form $N$ of degree $d$ in $d^2$ variables such that
$N(\mathbf z)=0$ for $\mathbf z\in K^{d^2}$ exactly when
$\mathbf z=0$.  Substitute
\[
                         (F_1,\ldots,F_s,0,\ldots,0)
\]
into this form.  Its zero condition is precisely the original
conjunction, and its degree is at most $dD$.
\end{proof}

For example, the preferred quartic presentations in 1315 and 1335
total coordinates come from 1272 and 1292 quadratic residuals,
respectively.  In both cases $\lceil\sqrt{s}\rceil=36$.
The number-field compiler therefore combines each residual system
into an equation of degree at most 72 without changing its
coordinate count.  The supplied determinant specifications have
degree exactly 72: the input occurs only as $+X_0^2$ in the
first residual, which multiplies the identity matrix, so the
coefficient of $X_0^{72}$ is one.  Over $\mathbb Q$ the
sum-of-squares quartic is better in degree.  Over a specified
nonreal number field a field-adapted reduced-norm construction preserves the
conjunction, whereas the sum of squares need not do so.  Neither
claim says that these base-changed residuals define the new ring
of integers.

\begin{corollary}\label{p3sys:minimum-variable-equality}
For a fixed subset $S$ of a number field $K$, the least total
number of coordinates in a definition of $S$ by one existential
polynomial equation equals the least total number in a definition
by a finite system of existential polynomial equations.  The
values are understood to be infinite if there is no such
definition.
\end{corollary}
\begin{proof}
A single equation is a finite system, giving one inequality.
The preceding combination preserves every coordinate, giving
the reverse inequality.
\end{proof}

Thus a system defining $\mathbb Z$ in at most ten total variables
would immediately give a single polynomial in at most ten total
variables.  If its equations are quadratic over $\mathbb Q$,
the combined polynomial would be quartic.  The system may be
much easier to discover or display, but it does not bypass a
lower bound on the variable count.

\subsection{Finite alternatives and nonzero conditions}

Suppose $P_j(T,\mathbf Y_j)=0$ are finitely many alternative
presentations, with $r_j$ witness coordinates.  Pad the lists to
$r=\max_jr_j$ using unused coordinates, and set
\[
                          P=\prod_j P_j.
\]
Over a field its existential projection is precisely the union
of the original projections.  Its total degree is at most
$\sum_j\deg P_j$, and no selector coordinate is needed.  This
also combines a finite union of systems after combining each
conjunction.  The assertion is about finite products with fixed
index sets, not a quantified integer choice.

Nonzero conditions require separate attention.  The conjunction
$G_1\ne0,\ldots,G_t\ne0$ is equivalent over a field to
\[
                         u\prod_{j=1}^tG_j-1=0
\]
with one additional witness $u$.  The new residual has degree at
most $1+\sum_j\deg G_j$.  This generic method therefore costs
one coordinate, even though conjunction and finite disjunction
of existing equations do not.  The cost cannot simply be ignored:
the one-coordinate set $K\setminus\{0\}$ is defined by the
nonzero condition $T\ne0$, but it is not the common zero set of
any finite family of one-variable polynomials.  Such a zero set
is finite unless every polynomial is zero, in which case it is
all of $K$.  Absorbing a guard into an existing witness requires
a separate equivalence, such as the reciprocal constructions
proved earlier.

\subsection{The degree of a universal combining polynomial}

Here a universal combining polynomial in $s$ inputs means a
polynomial $A\in K[Z_1,\ldots,Z_s]$ with
\[
              A(\mathbf z)=0\quad\Longleftrightarrow\quad
                           \mathbf z=0\qquad(\mathbf z\in K^s).
\]
Its defining property concerns all residual-value tuples.  It
is stronger than the property needed for one particular list
of residuals, whose value tuples may satisfy additional relations.

\begin{proposition}\label{p3sys:quadratic-combiner-optimality}
For $s\ge2$, a universal combining polynomial has degree at
least two.  If its degree is at most two, it is a homogeneous
anisotropic quadratic form.  Consequently:
\begin{enumerate}
\item if $K$ has a real embedding, the minimum degree is two
for every $s\ge2$;
\item if $K$ is totally imaginary, the minimum degree is two
for $2\le s\le4$, and three for $5\le s\le9$.
\end{enumerate}
For totally imaginary $K$ and $s\ge10$, the present argument
only gives
\[
 3\le\text{minimum universal-combiner degree}
                         \le\lceil\sqrt{s}\rceil.
\]
No optimality of the last upper bound is asserted.
\end{proposition}
\begin{proof}
The constant term is zero.  A nonzero linear form in at least
two variables has a nonzero vector in its kernel, so degree one
cannot work.  Write a polynomial of degree at most two as
$A=q+\ell$, with $q$ quadratic homogeneous and $\ell$ linear.
If both parts are nonzero, choose $v\in K^s$ with
$q(v)\ell(v)\ne0$.  The field is infinite, so two proper
hypersurfaces cannot cover its affine space.  With
$t=-\ell(v)/q(v)\ne0$ one has
\[
                          A(tv)=t^2q(v)+t\ell(v)=0,
\]
contradicting universality.  If $q=0$, the linear argument
already applies.  Therefore $\ell=0$, and $q$ must be
anisotropic.

A real embedding permits the anisotropic form
$Z_1^2+\cdots+Z_s^2$ in every dimension.  For a totally
imaginary number field, the $u$-invariant is four: every
quadratic form in at least five variables is isotropic, while
a quaternion division algebra supplies an anisotropic norm
form in four variables.  The upper restriction follows from
Hasse--Minkowski and the local $u$-invariant four; quaternion
existence follows by ramifying at two finite places.  See
\cite[Theorems~14.6.1 and 14.6.9]{Voight2021} and the explicit
summary in \cite[Section~5.1]{p3sysMerkurjevParimala2025}.
Restricting an anisotropic four-variable form to a coordinate
subspace treats $s=2,3,4$.  For $5\le s\le9$, a degree-three
division norm supplies the upper bound, and the quadratic
obstruction makes it sharp.  The general upper bound is
Corollary~\ref{p3sys:division-combination}.
\end{proof}

For $s=1$ the minimum is one.  These combining-polynomial
restrictions do not impose the same degree bound on a
particular definition: dependencies among its residuals can
allow cheaper combinations.  Conversely, the general degree
lower bound for a system defining $\mathcal O_K$ is only two,
since linear systems are affine spaces and have empty, constant,
or surjective input projection.  The much stronger restrictions
for quadratic systems were proved in
Section~\ref{p3sys:lower-bounds}.
